%% kapa: exam=1 round=37 year=2026 period=1 form=A subject=economics subject_ko=경제학원론 seq=10 number=50
%% source: Q-Net 감정평가사 기출문제 (cst00309) artlSeq=5258924 file=raw/1008/5258924/2026년 제37회 감정평가사 1차 1교시.pdf
%% conversion: kordoc 4.24.51 -> tools/gen1.py (auto); review=proofread-sonnet; flags=none
\begin{question}{50}
노동시장이 수요독점일 때, 이에 관한 설명으로 옳은 것은? (단, 생산물 시장은 완전경쟁이며, 노동수요곡선은 우하향, 노동공급곡선은 우상향한다.)
\begin{choices}
  \item 수요독점기업은 노동공급곡선 상에서 임금수준을 결정한다.
  \item 균형고용량은 완전경쟁 노동시장에 비해 수요독점 하에서 더 크다.
  \item 수요독점기업의 균형임금은 완전경쟁 노동시장의 균형임금보다 높다.
  \item 한계요소비용곡선(MFC)은 노동공급곡선보다 아래에 위치한다.
  \item 수요독점기업은 노동수요곡선을 가지고 있다.
\end{choices}
\end{question}
\answer{1}
